\documentclass[11pt]{article}
\usepackage[margin=1in]{geometry}
\usepackage{amsmath,amssymb,amsthm}
\usepackage[T1]{fontenc}
\usepackage{lmodern}
\usepackage{graphicx}
\usepackage[hidelinks]{hyperref}

\newtheorem{theorem}{Theorem}
\newtheorem{lemma}{Lemma}
\newtheorem{corollary}{Corollary}
\newtheorem{remark}{Remark}

\newcommand{\N}{\mathbb N}
\newcommand{\R}{\mathbb R}
\newcommand{\cI}{\mathcal I}
\newcommand{\Fin}{\mathrm{Fin}}
\newcommand{\Lac}{\mathcal{L}ac}
\newcommand{\Id}{\mathcal I_d}
\newcommand{\Prob}{\mathbb P}
\newcommand{\E}{\mathbb E}

\title{PLI is necessary below density zero:\\
a counterexample to the wPLI extension}
\author{Codex agent lane 16}
\date{July 19, 2026}

\begin{document}
\maketitle

\section*{Status and Source Question}

\noindent\textbf{Status:} counterexample likely valid.

Balcerzak--Poplawski--Wachowicz study, for an ideal $\cI$ on $\N$ and a real
series $\sum_n x_n$, the set
\[
 A(\cI,(x_n))
 =\left\{t\in\{0,1\}^{\N}:\sum_n t(n)x_n
        \text{ is $\cI$-convergent}\right\},
\]
equipped with the fair Bernoulli product measure $\lambda$ \cite{BPW}.  Their
Question 4.6 asks:

\begin{quote}
Does the statement of Theorem 4.3 remain true if PLI is replaced by wPLI?
In particular, does it work for the ideal from Example 3?
\end{quote}

\begin{center}
\includegraphics[width=0.95\linewidth]{figures/open_problem_crop.png}
\end{center}

The answer is negative even for the intended core assertion of Theorem 4.3,
namely the existence of a divergent real series with
$\lambda(A(\cI,(x_n)))=1$.  We prove a converse to the source theorem for all
ideals contained in the density-zero ideal.

\section*{Definitions}

For a real sequence $(z_n)$ and an ideal $\cI$ on $\N$, write
$z_n\to_{\cI}z$ when
\[
 \{n:|z_n-z|>\varepsilon\}\in\cI
 \qquad(\varepsilon>0).
\]
A series is $\cI$-convergent when its sequence of partial sums is
$\cI$-convergent.

The asymptotic-density-zero ideal is
\[
 \Id=\left\{B\subseteq\N:
       \lim_{N\to\infty}\frac{|B\cap[1,N]|}{N}=0\right\}.
\]
Thus $\Id$-convergence is statistical convergence.

An ideal $\cI$ has the \emph{property of long intervals} (PLI) if some
$E\in\cI$ contains intervals of every finite length.  It has the weak property
of long intervals (wPLI) if, for every fixed length $r$, some $E_r\in\cI$
contains infinitely many pairwise disjoint intervals of length $r$.

The source's ideal $\Lac$ is the smallest ideal containing every range
\[
 \{b_j:j\in\N\},\qquad b_1<b_2<\cdots,qquad b_{j+1}-b_j\longrightarrow\infty.
\]
Equivalently, its members are subsets of finite unions of such lacunary
ranges.  The source proves that $\Lac$ is dense, hence wPLI, but is not PLI
\cite[Remark 4.4(ii) and Example 3]{BPW}.

\section*{Proof Intuition}

The missing obstruction is that measure-one random convergence carries much
more information than convergence of the original series.  Pairing a selector
with its complement first shows that the original partial sums are
$\cI$-convergent.  When $\cI\subseteq\Id$, both the original partial sums and
almost every Bernoulli partial-sum path are statistically convergent.

Statistical convergence lets us choose deterministic endpoints at which the
random paths converge in probability.  After subtracting half of the original
partial sum, the increments become independent symmetric variables
$\pm x_n/2$.  Levy's maximal inequality says that convergence at the selected
endpoints controls every partial sum inside the intervening blocks.  A single
coefficient bounded away from zero would force a fixed-size excursion in its
block for every outcome.  Hence $x_n\to0$.

Now ordinary divergence becomes incompatible with ideal convergence unless
the ideal has PLI.  Indeed, a partial sum a fixed distance from the ideal limit
must, because the terms are small, remain far away for an arbitrarily long
interval of preceding indices.  One exceptional set in $\cI$ therefore
contains intervals of all lengths.

\section*{Formal Argument}

\begin{lemma}[Finite Levy maximal inequality]
\label{lem:levy}
Let $Y_1,\dots,Y_m$ be independent symmetric real random variables and put
$S_k=\sum_{i=1}^kY_i$.  Then, for every $a>0$,
\[
 \Prob\!\left(\max_{1\leq k\leq m}|S_k|>a\right)
 \leq 2\Prob(|S_m|>a).
\]
\end{lemma}

\begin{proof}
This is the standard reflection argument.  On the event that $S_k$ first
exits $[-a,a]$, write the remaining sum as $Z$.  Conditionally on the stopped
path, $Z$ and $-Z$ have the same law, while the triangle inequality shows
that at least one of $|S_k+Z|$ and $|S_k-Z|$ is greater than $a$.  Thus the
conditional probability that the terminal sum has absolute value greater
than $a$ is at least $1/2$.  Summing over the possible first exit times gives
the result.
\end{proof}

\begin{lemma}[Almost-sure statistical convergence forces null coefficients]
\label{lem:null}
Let $\cI\subseteq\Id$ and let $(x_n)$ be a real sequence.  If
\[
 \lambda(A(\cI,(x_n)))=1,
\]
then $x_n\to0$.
\end{lemma}

\begin{proof}
Write $s_n=\sum_{i=1}^n x_i$ and
$T_n(t)=\sum_{i=1}^n t(i)x_i$.  Let $A=A(\cI,(x_n))$ and let
$\phi(t)=\boldsymbol{1}-t$.  The map $\phi$ preserves $\lambda$, so
$A\cap\phi(A)$ has measure one and is nonempty.  Choose $t$ in this
intersection.  Both $\sum_i t(i)x_i$ and $\sum_i(1-t(i))x_i$ are
$\cI$-convergent.  Adding their partial sums shows that
\[
 s_n\longrightarrow_{\cI}s
\]
for some $s\in\R$.  Since $\cI\subseteq\Id$, $(s_n)$ converges statistically
to $s$.

For almost every selector $t$, the sequence $T_n(t)$ is likewise
statistically convergent; denote its statistical limit by $L(t)$.  The limit
may be chosen measurable.  For example, statistical convergence of $T_n(t)$
implies ordinary convergence of the Cesaro means of the bounded sequence
$\arctan(T_n(t))$ to $\arctan(L(t))$, which gives a measurable formula for
$L$ on the full-measure convergence set.

Fix $j\geq1$.  Statistical convergence and dominated convergence give
\[
 \frac1N\sum_{n=1}^N
 \lambda\bigl(|T_n-L|>1/j\bigr)\longrightarrow0.
 \tag{1}
\]
Also, the set
\[
 G_j=\{n:|s_n-s|\leq1/j\}
\]
has asymptotic density one.  Consequently, there are arbitrarily large
$n\in G_j$ for which
\[
 \lambda\bigl(|T_n-L|>1/j\bigr)<1/j;
\]
otherwise (1) would have positive lower limit.  Choose such integers
 $n_1<n_2<\cdots$, and put $n_0=0$.  Then
\[
 T_{n_j}\longrightarrow L\quad\text{in probability},
 \qquad s_{n_j}\longrightarrow s.
\]

Center the Bernoulli sums:
\[
 R_n=T_n-\frac12s_n
     =\sum_{i=1}^n\left(t(i)-\frac12\right)x_i.
\]
The summands are independent and symmetric, taking the values
$\pm x_i/2$.  We have
\[
 R_{n_j}\longrightarrow L-\frac12s
 \quad\text{in probability}.
\]
Thus the block increments
\[
 Z_j=R_{n_j}-R_{n_{j-1}}
\]
converge to zero in probability.  Applying
Lemma~\ref{lem:levy} to the $j$th block yields, for every $a>0$,
\[
 \Prob\!\left(
  \max_{n_{j-1}<k\leq n_j}|R_k-R_{n_{j-1}}|>a
 \right)
 \leq2\Prob(|Z_j|>a)\longrightarrow0.
 \tag{2}
\]

If some $i\in(n_{j-1},n_j]$ satisfied $|x_i|>4a$, then the corresponding
centered increment would have deterministic magnitude $|x_i|/2>2a$.
The partial sums immediately before and after that increment cannot both
have absolute value at most $a$.  Therefore the event on the left of (2)
would have probability one.  For all sufficiently large $j$ this contradicts
(2).  Hence
\[
 \max_{n_{j-1}<i\leq n_j}|x_i|\leq4a
\]
eventually.  Since $a>0$ was arbitrary and the blocks cover a tail of $\N$,
$x_n\to0$.
\end{proof}

\begin{lemma}[Divergent ideal convergence with null terms forces PLI]
\label{lem:pli}
Suppose the real series $\sum_n x_n$ is $\cI$-convergent but not ordinarily
convergent, and $x_n\to0$.  Then $\cI$ has PLI.
\end{lemma}

\begin{proof}
Let $s_n=\sum_{i=1}^n x_i$ and $s_n\to_{\cI}s$.  Since $(s_n)$ does not
converge ordinarily to $s$, there are $\delta>0$ and arbitrarily large $n$
such that $|s_n-s|>\delta$.  Put
\[
 E=\{m:|s_m-s|>\delta/2\}\in\cI.
\]
Fix a positive integer $r$.  Choose one of the above indices $n$ so large
that $n\geq r$ and
\[
 |x_i|<\frac{\delta}{2r}\qquad(i\geq n-r+1).
\]
For $m\in\{n-r+1,\dots,n\}$,
\[
 |s_m-s|
 \geq |s_n-s|-\sum_{i=m+1}^n|x_i|
 >\delta-\frac{r-1}{2r}\delta>\frac\delta2.
\]
Thus $E$ contains an interval of length $r$.  As $r$ was arbitrary, $E$ is a
PLI witness.
\end{proof}

\begin{theorem}[PLI necessity below density zero]
\label{thm:necessity}
Let $\cI\subseteq\Id$.  If there is a divergent real series $\sum_nx_n$ such
that
\[
 \lambda(A(\cI,(x_n)))=1,
\]
then $\cI$ has PLI.
\end{theorem}

\begin{proof}
By the complement-selector argument in Lemma~\ref{lem:null}, the original
series is $\cI$-convergent.  The same lemma gives $x_n\to0$.
Lemma~\ref{lem:pli} now yields PLI.
\end{proof}

\begin{corollary}[Classification below density zero]
\label{cor:classification}
Let $\cI\subseteq\Id$.  The following are equivalent:
\begin{enumerate}
 \item $\cI$ has PLI;
 \item there is a divergent real series $\sum_nx_n$ for which
       $\lambda(A(\cI,(x_n)))=1$.
\end{enumerate}
\end{corollary}

\begin{proof}
Theorem~\ref{thm:necessity} proves $(2)\Rightarrow(1)$.  The implication
$(1)\Rightarrow(2)$ is the core conclusion of source Theorem 4.3; its explicit
block construction uses $2^k$ terms $2^{-k}$ followed by $2^k$ terms
$-2^{-k}$ inside each PLI interval and applies Chebyshev, Borel--Cantelli, and
the Cauchy criterion \cite[Theorem 4.3]{BPW}.
\end{proof}

\section*{The wPLI Counterexample}

\begin{lemma}
\label{lem:lac-density}
$\Lac\subseteq\Id$.
\end{lemma}

\begin{proof}
Let $B=\{b_1<b_2<\cdots\}$ with $b_{j+1}-b_j\to\infty$.  Given $q\geq1$,
all sufficiently late gaps in $B$ are at least $q$.  Hence
\[
 |B\cap[1,N]|\leq C_q+\frac{N}{q}
\]
for a constant $C_q$.  Its upper asymptotic density is at most $1/q$, and
therefore is zero.  Finite unions of such ranges also have density zero, as
do their subsets.  These are exactly the members of $\Lac$.
\end{proof}

\begin{theorem}[Negative answer to Question 4.6]
\label{thm:counterexample}
The wPLI hypothesis does not suffice in Theorem 4.3 of the source paper.  More
precisely, the ideal $\Lac$ from source Example 3 has wPLI, but there is no
divergent real series $\sum_nx_n$ satisfying
\[
 \lambda(A(\Lac,(x_n)))=1.
\]
\end{theorem}

\begin{proof}
The source proves that $\Lac$ is dense and consequently has wPLI, while its
Example 3 proves that $\Lac$ does not have PLI.  By
Lemma~\ref{lem:lac-density}, $\Lac\subseteq\Id$.  If the displayed series
existed, Theorem~\ref{thm:necessity} would force $\Lac$ to have PLI, a
contradiction.
\end{proof}

\section*{Printed-Theorem Erratum}

Theorem 4.3 additionally prints the conclusion $x_n\not\to0$.  The displayed
construction in its proof actually uses coefficients $\pm2^{-k}$ on the
$k$th active block, so its coefficients tend to zero.  This is not merely a
missing observation: Lemma~\ref{lem:null} proves that whenever
$\cI\subseteq\Id$ and $\lambda(A(\cI,(x_n)))=1$, one necessarily has
$x_n\to0$.  In particular, the extra printed conclusion is impossible for
$\cI=\Id$, even though $\Id$ has PLI.  The core measure-one construction used
in Corollary~\ref{cor:classification} is unaffected.

The same mismatch appears in the source's concluding claim that its constructed
series has $\sum_nx_n^2=\infty$.  In fact the $k$th active block contributes
\[
 2^{k+1}(2^{-k})^2=2^{1-k},
\]
so the sum of squares is finite.  These source issues are logically separate
from Theorem~\ref{thm:counterexample}, which disproves the intended wPLI
extension itself.

\section*{Verification and Novelty Check}

The proof is non-computational.  Its only probabilistic input,
Lemma~\ref{lem:levy}, is included with the standard reflection proof.  The
two deterministic checks are explicit: every $\Lac$ member has density zero,
and a divergent ideal-convergent null-term series creates arbitrarily long
intervals in one exceptional set.

The run indexes were searched for arXiv:1803.03699, ideal-convergent
subseries, PLI, wPLI, the lacunary ideal, and measure-one Bernoulli subseries;
no duplicate was found.  A bounded web search on July 19, 2026 used the exact
question, the exact paper title, the DOI, and the phrases ``weak property of
long intervals'' and ``property of long intervals.''  Semantic Scholar listed
six citing records.  The accessible citing sources---including
\emph{Another characterization of meager ideals}, \emph{Convergent subseries
of divergent series}, and \emph{Ideal boundedness of subseries and
rearrangements in Banach spaces vs Banach spaces possessing a copy of
$c_0$}---cite the source only for adjacent category, summable-ideal, or
boundedness results; none states or proves an answer to Question 4.6.  No
later answer was found.  This supports, but cannot certify, novelty.

Human-review recommendation: check first the selection of $(n_j)$ in
Lemma~\ref{lem:null} and the blockwise application of Levy's inequality.
Those are the only steps where measure and statistical convergence interact.

\begin{thebibliography}{9}

\bibitem{BPW}
M. Balcerzak, M. Poplawski, and A. Wachowicz,
\emph{Ideal convergent subseries in Banach spaces},
Quaestiones Mathematicae 42 (2019), no. 6, 765--779;
arXiv:1803.03699.\newline
DOI: \href{https://doi.org/10.2989/16073606.2018.1497724}
{10.2989/16073606.2018.1497724}.

\end{thebibliography}

\end{document}
